\section{Passing between real and complex spaces}\label{sec:fields}


Theorem~\ref{pos:theorem} has established the positive implication for
real Banach spaces. We now transfer it to complex spaces. We also prepare
the transfer needed for the negative implication: after constructing a
real counterexample, we will take its complexification. The two transfers
require different arguments.

\begin{proposition}\label{ev:positive-transfer}
Let $X$ be an infinite-dimensional complex Banach space, and let $X_{\R}$
be its underlying real space.  If $X_{\R}$ has a separable
infinite-dimensional real quotient, then $X$ has a separable
infinite-dimensional complex quotient.
\end{proposition}

\begin{proof}
The map
\[
 A:(X_{\R})^*\longrightarrow (X^*)_{\R},\qquad
 (Ag)(x)=g(x)-i g(ix),
\]
is a real-linear isometry with inverse $f\mapsto\operatorname{Re}f$.
Indeed $Ag$ is complex-linear, and phase rotation in $B_X$ proves
$\|Ag\|=\|g\|$.
Apply Proposition~\ref{ev:criterion} to $X_{\R}$, and transport its
witness under $A$ to a real closed infinite-dimensional subspace
$E\subset (X^*)_{\R}$.  The complex-valued real-linear evaluations on
$E$ have norm-separable image: under $A$, evaluation at $x$ is the
difference of real evaluation at $x$ and $i$ times real evaluation at
$ix$.  Evaluation on $iE$ also has norm-separable image.

Consider the real operator
\[
 q:iE\longrightarrow (X^*)_{\R}/E,\qquad q(w)=w+E.
\]
First suppose that $q$ is not strictly singular.  There are a closed
infinite-dimensional $W\subset iE$ and $c>0$ such that
$\operatorname{dist}(w,E)\ge c\|w\|$ for $w\in W$.
Since $iW\subset E$, the space $H=W+iW$ is a complex subspace whose
displayed sum is direct over $\R$, and
\[
 \|w\|\le c^{-1}\|w+iv\|,\qquad
 \|v\|\le(1+c^{-1})\|w+iv\|\quad(w,v\in W).
\]
These estimates make $H$ closed and its real coordinate projections
bounded. The two coordinate projections combine the norm-separable
evaluation images on $W$ and $iW$ into a separable set of real-linear
maps $H\to\C$. The original evaluations are complex linear, with
the same norm, so $R_H(X)$ is norm separable in $H^*$.
Proposition~\ref{ev:criterion}, over
$\C$, gives the desired quotient.

Now suppose that $q$ is strictly singular.  Fix positive numbers
$\varepsilon_n$ with $\sum_n\varepsilon_n<\infty$.
Choose $f_n\in E$ of norm one inductively, requiring that all later
vectors vanish on finite actual $1/2$-norming sets $D_m\subset B_X$ for
the \emph{complex} spans of the earlier vectors.  Each such vanishing
condition consists of two real linear equations, so their common kernel
in $E$, say $M$, is closed and infinite dimensional. Thus $iM$ is
a closed infinite-dimensional subspace of $iE$. Strict singularity of
$q|_{iM}$ lets us choose $f_n$ with
$\operatorname{dist}(if_n,E)<\varepsilon_n/2$.
Choose $g_n\in E$ with $\|g_n-if_n\|<\varepsilon_n$.
Lemma~\ref{ev:prefix} shows that $(f_n)$ is complex basic, with
coefficient bound $4$.

Put $H=\overline{\operatorname{span}}_{\C}\{f_n:n\ge1\}$.
For finite sums, define a real-linear map by
\[
 L\left(\sum_n(a_n+ib_n)f_n\right)
       =\sum_n(a_nf_n+b_ng_n),\qquad a_n,b_n\in\R.
\]
The operator $K=L-I$ extends to $H$, because
\[
 Kh=\sum_n\operatorname{Im}\bigl(f_n^\#(h)\bigr)(g_n-if_n),
 \qquad
 \|K-K_N\|\le4\sum_{n>N}\varepsilon_n,
\]
where $K_N$ is the finite partial sum.  In particular, $K$ is a norm
limit of finite-rank real-linear operators.  Since $E$ is closed and
$L$ maps finite sums into $E$, we have $L(H)\subset E$.

For $x\in X$, write $e_x(h)=h(x)$ on $H$ and also use $e_x$ for
evaluation on $X^*$.  Work temporarily in the normed space of bounded
real-linear maps from $H$ to $\C$.
The family $\{e_x\circ L:x\in X\}$ is norm separable, by the evaluation
property of $E$ and boundedness of $L$.
For $x\in B_X$, the functionals $e_x\circ K_N$ form a bounded subset of
a finite-dimensional space, and
\[
 \sup_{x\in B_X}\|e_x\circ K-e_x\circ K_N\|
 \le\|K-K_N\|\longrightarrow0.
\]
Thus $\{e_x\circ K:x\in B_X\}$ is totally bounded.  The identity
$e_x=e_x\circ L-e_x\circ K$ proves norm separability of the evaluations
on $B_X$, and their countably many dilates give the same conclusion on
$X$.  The complex-linear functionals on $H$ embed isometrically in the
space of real-linear maps just used.  Hence $R_H(X)\subset H^*$ is norm
separable, and Proposition~\ref{ev:criterion} applies again.
\end{proof}

Apply Theorem~\ref{pos:theorem} to the underlying real space of a complex
Banach space, then apply Proposition~\ref{ev:positive-transfer}. This
proves the positive implication over $\C$. To transfer a counterexample, we first show that
a finite real direct sum cannot acquire a separable infinite-dimensional
quotient when neither summand has one.

\begin{lemma}\label{ev:finite-sum}
If neither of two real Banach spaces $X_1,X_2$ has a separable
infinite-dimensional quotient, then their direct sum has none.
\end{lemma}

\begin{proof}
Use the norm $\|(x_1,x_2)\|=\max(\|x_1\|,\|x_2\|)$ on
$Y=X_1\oplus X_2$, so $Y^*=X_1^*\oplus_1X_2^*$.
If $Y$ had such a quotient, Proposition~\ref{ev:criterion} would provide
a closed infinite-dimensional $F\subset Y^*$ with norm-separable
evaluation image.  Let $p_j:F\to X_j^*$ be its coordinate operators.

If $p_1$ is not strictly singular, choose a closed infinite-dimensional
$G\subset F$ on which it is bounded below.  Otherwise,
Lemma~\ref{ev:small} gives a closed infinite-dimensional $G\subset F$
with $\|p_1|_G\|\le1/2$, and then
\[
 \|p_2g\|=\|g\|-\|p_1g\|\ge\tfrac12\|g\|\qquad(g\in G).
\]
Thus some $p_j$ identifies $G$ with a closed infinite-dimensional
$H\subset X_j^*$.  For $x\in X_j$, evaluation on $H$ is obtained by
evaluating the corresponding coordinate vector of $Y$ on $F$,
restricting to $G$, and composing with $(p_j|_G)^{-1}$.
Its image is norm separable.  Proposition~\ref{ev:criterion} gives the
forbidden quotient of $X_j$.
\end{proof}

\begin{corollary}\label{ev:fields}
The universal separable quotient assertions over $\R$ and over $\C$
are equivalent.  More specifically, if a real Banach space $X$ has no
separable infinite-dimensional quotient, its complexification has no
separable infinite-dimensional complex quotient.
\end{corollary}

\begin{proof}
Proposition~\ref{ev:positive-transfer} proves the positive implication
from real to complex spaces.  For the other implication, suppose $X$
is a real counterexample.  Give $X_{\C}=X+iX$ the norm
\[
 \|x+iy\|_{\C}
   =\sup_{\theta\in\R}\|\cos\theta\,x-\sin\theta\,y\|.
\]
Phase rotation proves absolute homogeneity over $\C$, and the triangle
inequality follows from that of $X$.  Moreover,
\[
 \max(\|x\|,\|y\|)\le\|x+iy\|_{\C}
 \le\|x\|+\|y\|\le2\max(\|x\|,\|y\|).
\]
Thus $X_{\C}$ is complete and its underlying real space is isomorphic to
$X\oplus_\infty X$.  A separable infinite-dimensional complex quotient
would give a separable infinite-dimensional real quotient of this sum,
contrary to Lemma~\ref{ev:finite-sum}.
\end{proof}
